\chapterhead{55}{ORR-16}{Case Study: Model Risk and Model Validation}{1}{3}

\lohead{55}{a}{Define a model and describe different ways that financial
institutions can become exposed to model risk.}

\orsub{Model risk exposure}

\begin{ordefbox}{Model risk}
Model risk refers to potential adverse outcomes from using models that apply
statistical, economic, financial, or mathematical theories and assumptions to
process data into estimates.
\end{ordefbox}

\orsubb{Types of model risk}

\noindent\begin{minipage}{\textwidth}
\renewcommand{\arraystretch}{1.35}
\begin{tabularx}{\textwidth}{@{}L{0.22\textwidth}Y@{}}
\rowcolor{navy} \thd{Type} & \thd{What goes wrong}\\
{\tabfont \textbf{Execution risk}} &
{\tabfont The risk of models not functioning as intended due to errors in input
data or in coding of the model.}\\
\rowcolor{paleblue}
{\tabfont \textbf{Conceptual errors}} &
{\tabfont Arise when the assumptions are invalid, meaning they do not represent
reality, or when incorrect modelling techniques are used.}\\
\end{tabularx}
\end{minipage}
\orfigcap{Execution risk is about the build; conceptual error is about the idea.}

% ---------------------------------------------------------------------
\lohead{55}{b}{Describe the role of the model risk management function and
explain best practices in the model risk management and validation processes.}

\begin{orkeybox}
{\sffamily\bfseries\fontsize{8.6}{10}\selectfont\color{navy}Model Risk Management
(MRM) team}\par
\vspace{2pt}
An independent team of experts tasked with mitigating model risk, setting
standards for documentation, data quality, and model control.
\end{orkeybox}

\orsub{What MRM does}

\begin{lettered}
\item \textbf{Risk tiers.} Not all models are equally risky. The MRM team assigns
models to different risk levels based on how important their outputs are, their
complexity, and how they are used. Higher risk models get more scrutiny.
\item \textbf{Validation.} MRM makes sure models are validated regularly. MRM is
an ongoing process, not just periodic reviews. This involves checking the model's
assumptions, testing how it performs with different data, and comparing its
outputs to real-world results.
\item \textbf{Monitoring.} MRM keeps an eye on models over time to make sure they
are still working as expected, especially as market conditions change.
\end{lettered}

% ---------------------------------------------------------------------
\lohead{55}{c}{Describe lessons learned from the three case studies involving
model risk.}

\orsub{Model risk case studies}

\orsubb{1) Gaussian copula and CDO pricing}

The Gaussian copula function was used to price collateralised debt obligations
(CDOs). This model assumed market efficiency and used credit default swap (CDS)
prices to infer correlations among assets in a collateral pool. The model failed
to adequately adjust to rapidly changing market conditions during the 2008
financial crisis, leading to significant mispricing of CDOs.

\begin{orexambox}{Lesson}
MRM would have identified limitations in the CDO pricing model and promoted
better communication.
\end{orexambox}

\orsubb{2) The Barclays acquisition of Lehman Brothers' assets, and a spreadsheet error}

\begin{lettered}
\item During the liquidation of Lehman Brothers, Barclays planned to acquire
specific assets and sent a spreadsheet to their lawyers, which included rows of
assets they did not wish to acquire, hidden.
\item A junior associate then converted this spreadsheet to PDF for court
submission without realising that the hidden rows became visible. This led to an
inadvertent claim on unwanted assets.
\end{lettered}

\begin{orexambox}{Lesson}
The importance of careful data management, and the potential consequences of
implementation errors in the use of software tools. MRM could have prevented the
Barclays spreadsheet error by advocating for clear data presentation.
\end{orexambox}

\orsubb{3) NASA Mars Orbiter}

The Mars Orbiter is an example of a rather innocuous mistake, which led to the
loss of a \$125 million satellite. In this case, the engineering team at Lockheed
Martin used English units of measurement (commonly used in the United States),
while NASA's convention was to use the metric system. This mistake in measurement
of model inputs led to the loss of a multi-million-dollar satellite.

\begin{orexambox}{Lesson}
MRM would have ensured proper unit conversions for the Mars Orbiter project.
\end{orexambox}
